% Bibliografía
\chapter*{Referencias}

\addcontentsline{toc}{chapter}{Referencias}
[1] M. H. Al-Saleh y U. Sundararaj, “Electromagnetic interference shielding mechanisms of CNT/polymer composites“, Carbon, vol. 47, no. 7, pp. 1738–1746, 2009.


[2] W. Bauhofer y J. Z. Kovacs, “A review and analysis of electrical percolation in carbon nanotube polymer composites“, Composites Science and Technology, vol. 69, no. 10, pp. 1486–1498, 2009.


[3] O. A. Castañeda-Uribe, “Polymer Nanocomposites: Synthesis and Correlations Between Distributions and Properties“, Tesis de Doctorado, Dept. Ing. Eléc. y Electrónica, Univ. de los Andes, Bogotá, Colombia, 2016.


[4] O. A. Castañeda-Uribe y A. Avila, “Enhancing Electromagnetic Interference Shielding Effectiveness of Polymer Nanocomposites by Modifying Subsurface Carbon Nanotube Distribution“, Advanced Engineering Materials, vol. 22, no. 9, p. 2000707, 2020.


[5] O. A. Castañeda-Uribe, R. Reifenberger, A. Raman y A. Avila, “Depth-Sensitive Subsurface Imaging of Polymer Nanocomposites Using Second Harmonic Kelvin Probe Force Microscopy“, ACS Nano, vol. 9, no. 3, pp. 2938–2947, 2015.


[6] R. C. Gonzalez y R. E. Woods, Digital Image Processing, 4ta ed. Nueva York, NY, EE. UU.: Pearson, 2017.


[7] T. Hastie, R. Tibshirani y J. Friedman, The Elements of Statistical Learning: Data Mining, Inference, and Prediction, 2da ed. Nueva York, NY, EE. UU.: Springer, 2009.


[8] P. Klapetek, Quantitative Data Processing in Scanning Probe Microscopy: SPM Applications for Nanometrology, 2da ed. Oxford, U.K.: Elsevier, 2018.


[9] N. Otsu, “A Threshold Selection Method from Gray-Level Histograms“, IEEE Transactions on Systems, Man, and Cybernetics, vol. 9, no. 1, pp. 62–66, 1979.


[10] D. R. Paul y L. M. Robeson, “Polymer nanotechnology: Nanocomposites“, Polymer, vol. 49, no. 15, pp. 3187–3204, 2008.


[11] B. W. Silverman, Density Estimation for Statistics and Data Analysis. Londres, U.K.: Chapman \& Hall, 1986.


[12] J.-M. Thomassin, C. Jérôme, T. Pardoen, C. Bailly, I. Huynen y C. Detrembleur, “Polymer/carbon based composites as electromagnetic interference (EMI) shielding materials“, Materials Science and Engineering: R: Reports, vol. 74, no. 7, pp. 211–232, 2013.
